\subsection{nml\_limiter}
\label{SEC:nml_limiter}
\inp*[nmlLimiter]{}[nml\_limiter][]

This \hyperref[SEC:User_Input]{namelist block}
is used to specify limiter, filter, and marker controls.

\begin{nmlblocklst}[
,label={LST:nml_limiter}
,caption={Namelist nml\_limiter.}
]
&nml_limiter
  Global_Filter_Interval      =  0
  Global_Filter_Option        =  0
  Global_Filter_Order         = 16
  Global_Filter_Alpha         = 36.0
  Global_Filter_Etac          =  0.0
  Sponge_Filter_Interval      =  0
  Sponge_Filter_Option        =  0
  Sponge_Filter_Order         =  8
  Sponge_Filter_Alpha         = 36.0
  Sponge_Filter_Etac          =  0.0
  Sponge_Filter_X_Start       = huge(0.0)  ! [input grid units]
  Local_Limiter_Interval      =  0
  Local_Limiter_Option        =  0
  Local_Filter_Order          = 16
  Local_Filter_Alpha          = 36.0
  Local_Filter_Etac           =  0.0
  Marker_Option               =  0
  Marker_Variable             =  0
  Marker_Min                  =  1.0
  Marker_Max                  =  1.0
  Marker_Dpsi                 = -1.0
  Marker_Stored_Interval      =  1
  Thermofix%Method            =  0
  Thermofix%Iterations        = 10
  Thermofix%Check_Cells       =  T
  Thermofix%Check_Faces       =  T
  Thermofix%Check_Nodes       =  T
  Thermofix%Verbose           =  T
  Thermofix%Scale_Factor      = 0.5
  Thermofix%Min_Pressure      = 10000.  ! [Pa]
  Thermofix%Min_Temperature   =   100.  ! [K]
  Thermofix%Min_Density       = 0.001   ! [kg/m^3]
  Positivity%Option           = 0
  Positivity%Density_Limit    = 1.0e-6  ! [kg/m^3]
  Positivity%Pressure_Limit   = 1.0e-6  ! [Pa]
  Entropy_Limiter%Option      = 0
  Entropy_Limiter%Interpolant = 0
  Entropy_Limiter%Tolerance   = 0.001
  Entropy_Limiter%Check_Faces = F
/
\end{nmlblocklst}

%==========================================================================

%\begin{nmldescitem}
%Pyramid_Filter_Option = 1
%\end{nmldescitem}
%\inp*[nmlLimiter::PyramidFilterOption]{}[Pyramid\_Filter\_Option][]
%
%\begin{nmldescbody}
%This input specifies how to filter the modes of pyramid elements.
%%
%All options use a tensor product of a 1D exponential filter with
%a 2D quad exponential filter to generate the pyramid filter matrices for
%all exponential filters: global, limiter, and sponge.
%
%The 1D filter is relative to the pyramid height from base to peak, with the
%base being the constant mode, and the peak being the \p{}-th mode, where
%$\p{}=\inp[nmlNumerics::SolutionOrder]{Solution\_Order}$.
%%
%The same 1D filter is used for all options below.
%
%The 2D filter at the base is equivalent for all options below, and is also
%the same as the 2D filter for an actual quad element of order \p{}.
%%
%Each increase in height of the pyramid subtracts 1 from the order of the
%quad used to create the 2D filter (\thatis the base is a quad of order \p{},
%the next level up is a quad of order \p{}-1, then a quad of order \p{}-2,
%and eventually at the peak a ``quad'' of order \p{}-\p{}=0).
%%
%The options below essentially amount to differences in how the 2D filter
%is defined for all of the pyramid layers above the base.
%
%Let $k$ be the level of the pyramid ranging from 0 to \p{}, where 0 is the
%base and \p{} is the peak.
%%
%Let the cutoff order \pc{} at level $k$, denoted by \pc{\p{}-k}{},
%be defined as
%  $\pc{\p{}-k}=\eta_c * \mathrm{real}(\p{}{-k})$,
%where $\eta_c$ is the cutoff ratio defined in
%\sref{SEC:Filtering}.
%
%The permissible values for \inptext{Pyramid\_Filter\_Option}
%can be summarized as:
%\\
%\begin{tabular}{ @{\qquad} r @{ = } p{0.85\linewidth} @{} }
%    1 & Scale $\eta_c$ so that the cutoff order \pc{} remains constant
%        relative to \p{} for all $k$. Note that $\eta_c$ will take the
%        value 1 if it evaluates to something larger than 1.\\
%    2 & Use the same value of $\eta_c$ for each level of the pyramid.
%        This means the cutoff order will change at each level.
%        For example, if
%        \p{}=3 and $\eta_c$=0.5: at the base level ($k$=0)
%        $\pc{\p{}-0}=0.5*\mathrm{real}(\p{}{-0})=2$ and at the second level
%        ($k$=2) $\pc{\p{}-2}=0.5*\mathrm{real}(\p{}{-2})=1$.
%        This results in a slightly more dissipative filter than the
%        first option.\\
%    3 & Use a full 3D tensor product of the 1D exponential filter.
%        This is equivalent to extracting only the pyramid modes from
%        the hex filter.
%\end{tabular}
%
%\NOTE
%After examining how the modes of other cell geometries are filtered,
%the default value of 1 appears to be the ``most correct'' way
%to filter the modes of pyramid cells.
%\end{nmldescbody}

%==========================================================================

\begin{nmldescitem}
Global_Filter_Interval = 0
\end{nmldescitem}
\inp*[nmlLimiter::GlobalFilterInterval]{}[Global\_Filter\_Interval][]

\begin{nmldescbody}
Specifies the frequency to apply the global filter.
\\
\begin{tabular}{ @{\quad} r @{ \quad } p{0.85\linewidth} @{} }
  $\le 0$ & Apply the global filter after each Runge-Kutta stage. \\
 %$> 0$   & Apply the global filter after the last Runge-Kutta stage
 %          of every
 %          $\inptext{Global\_Filter\_Interval}^{\;\textrm{th}}$
 %          time step. \\
  $> 0  $ & Apply the global filter after the last Runge-Kutta stage if
            \newline
            $\textrm{mod} \left(         {current\_iteration},
                                 \inptext{Global\_Filter\_Interval}
                          \right) = 0$.
            \newline
            The ${current\_iteration}$ is the total aggregate iteration
            number including previous jobs.
 %          \\
 %$< 0  $ & The global filter will NOT be applied even if
 %          $\inptext{Global\_Filter\_Option} \neq 0$
\end{tabular}
\end{nmldescbody}

%==========================================================================

\begin{nmldescitem}
Global_Filter_Option = 0
\end{nmldescitem}
\inp*[nmlLimiter::GlobalFilterOption]{}[Global\_Filter\_Option][]

\begin{nmldescbody}
Specifies whether the global filter is used for stabilization.
%
This filter is denoted as a global filter because it is used on
all cells regardless of the quality of the solution in each cell.
%
The available options are as follows:
\\
\begin{tabular}{ @{\quad} r @{\quad} p{0.80\linewidth} @{} }
  $0=$ & Disable the global filter.\\
  $1=$ & Use the exponential filter, \eref{EQN:ExpFilterMatrix}, 
         where the polynomial order is given by 
         $\PO=\inp[nmlNumerics::SolutionOrder]{Solution\_Order}$
         for 1D, 2D, and 3D.
\end{tabular}

GFR uses a three-pass filtering approach involving global, local, and
sponge filters.
%
Each of them and their sequence of application is described in
\sref{SEC:Filtering}.
%
Generally, we recommend that the global filter not be used, because it can
introduce significant dissipation.
%
When it is absolutely necessary to use the global filter, the filter settings
should be set as weak as possible.
\end{nmldescbody}

%==========================================================================

\begin{nmldescitem}
Global_Filter_Order = 16
\end{nmldescitem}
\inp*[nmlLimiter::GlobalFilterOrder]{}[Global\_Filter\_Order][]

\begin{nmldescbody}
This is an integer value between ($2 \leq s \leq 32$) that specifies the
exponent\slash order of the global exponential filter,
\eref{EQN:ExpFilterMatrix}.
%
As discussed in \sref{SEC:Filtering} and shown in \fref{FIG:Filter_s},
the filter order controls how rapidly the filter decays.
Larger values act to concentrate the filtering toward the highest mode.

Because the global filter is applied to all cells regardless of the
solution quality in each cell, one generally wants to use
as large a value as possible so as not to affect the lower modes.
%
Starting with \inptext{Global\_Filter\_Order} of 32 or 16
and halving if necessary is a good approach.
\end{nmldescbody}

%==========================================================================

\begin{nmldescitem}
Global_Filter_Alpha = 36.0
\end{nmldescitem}
\inp*[nmlLimiter::GlobalFilterAlpha]{}[Global\_Filter\_Alpha][]

\begin{nmldescbody}
This is a real value between ($0. < \alpha \leq 36.$) that specifies the damping
parameter of the global exponential filter, \eref{EQN:ExpFilterMatrix}.
%
As discussed in \sref{SEC:Filtering}, and shown in \fref{FIG:Filter_alpha},
this parameter controls the level of damping applied to the highest mode;
larger values apply more damping.
%
A value of 36.0 effectively filters the highest mode to machine zero.

Because the global filter is applied to all cells regardless of the
solution quality in each cell, one ideally wants to use
as small a value as possible so as not to introduce unwanted dissipation.
%
However, since numerical error often first appears in the highest modes
it may be necessary to aggressively filter those modes before it
contaminates the lower modes.
%
Starting with \inptext{Global\_Filter\_Alpha} of 2 or 4
and doubling if necessary is a good approach.

In the special case of \PO[=1], the \PO[1] mode will be multiplied by
$e^{-\alpha}$.
\end{nmldescbody}

%==========================================================================

\begin{nmldescitem}
Global_Filter_Etac = 0.0
\end{nmldescitem}
\inp*[nmlLimiter::GlobalFilterEtac]{}[Global\_Filter\_Etac][]

\begin{nmldescbody}
This is a real value between ($0. \leq \eta_c \le 1.$) that specifies the modal
cutoff fraction for the global exponential filter, \eref{EQN:ExpFilterMatrix}.
%
As discussed in \sref{SEC:Filtering} and shown in \fref{FIG:Filter_etac},
modes at or below this cutoff fraction will be entirely preserved.
%
A value of $1.0$ applies no filter and wastes computational resources.

For most simulations it is desirable to preserve the linear mode,
and so a good approach is to set this input value to $\eta_c=1/{\PO}$ where
$\PO=\inp[nmlNumerics::SolutionOrder]{Solution\_Order}$.
%
However, if ${\PO} \leq 2$ it may be best to use a value of 0,
and if ${\PO} > 4$ it may be possible to use a larger fraction.

In the special case of \PO[=1], any value $<1.0$ will just multiply the \PO[1]
mode by $e^{-\alpha}$.
\end{nmldescbody}

%==========================================================================

\begin{nmldescitem}
Sponge_Filter_Interval =  0
Sponge_Filter_Option   =  0
Sponge_Filter_Order    =  8
Sponge_Filter_Alpha    = 36.0
Sponge_Filter_Etac     =  0.0
Sponge_Filter_X_Start  = huge(0.0)  ! [input grid units]
\end{nmldescitem}
\inp*[nmlLimiter::SpongeFilterInterval]{}[Sponge\_Filter\_Interval][]
\inp*[nmlLimiter::SpongeFilterOption]{}[Sponge\_Filter\_Option][]
\inp*[nmlLimiter::SpongeFilterOrder]{}[Sponge\_Filter\_Order][]
\inp*[nmlLimiter::SpongeFilterAlpha]{}[Sponge\_Filter\_Alpha][]
\inp*[nmlLimiter::SpongeFilterEtac]{}[Sponge\_Filter\_Etac][]
\inp*[nmlLimiter::SpongeFilterXStart]{}[Sponge\_Filter\_X\_Start][]

\begin{nmldescbody}
GFR uses a three-pass filtering approach involving global, local, and
sponge filters.
%
Each of them and their sequence of application is described in
\sref{SEC:Filtering}.
%
Setting \inptext{Sponge\_Filter\_Option}=1 activates the sponge filter,
which is ideally used as an extra-strong exponential filter to damp
non-physical pressure reflections near the outflow boundary.
%
It is applied to any cells where all the $x$-coordinate values
(in \key[key:GridUnits]{grid units}) are larger than
that specified by \inptext{Sponge\_Filter\_X\_Start}.

The details for these inputs are the same as for the similarly
named \inptext{Global\_Filter} options.
% 
The user must manually specify the starting $x$-location
(in \key[key:GridUnits]{grid units}) of the sponge region via
\\% Prevent next word from entering right margin.
\inptext{Sponge\_Filter\_X\_Start},
since the default value of \inptext{huge(0.0)} will be beyond
any current grid location.

\noindent\NOTE
If $\inp[nmlReference::AoAAlpha]{AoA\_Alpha}=\pm90$\textdegree{} and
$\inp[nmlReference::AoABeta]{AoA\_Beta}=0$\textdegree{}, the $y$-axis will be
assumed as the axial coordinate direction, and this input will be applied to the
$y$-coordinate.
%
If $\inptext{AoA\_Alpha}=0$\textdegree{} and
$\inptext{AoA\_Beta}=\pm90$\textdegree{}, the $z$-axis will be assumed as the
axial coordinate direction, and this input will be applied to the
$z$-coordinate.
\end{nmldescbody}

%==========================================================================

\FINALFIX{\newpage}
\begin{nmldescitem}
Local_Limiter_Interval = 0
\end{nmldescitem}
\inp*[nmlLimiter::LocalLimiterInterval]{}[Local\_Limiter\_Interval][]

\begin{nmldescbody}
Specifies the frequency to apply the local limiter.
\\
\begin{tabular}{ @{\quad} r @{ \quad } p{0.85\linewidth} @{} }
  $\le 0$ & Apply the limiter after each Runge-Kutta stage. \\
  $> 0  $ & Apply the limiter after the last Runge-Kutta stage if
            \newline
            $\textrm{mod} \left(         {current\_iteration},
                                 \inptext{Local\_Limiter\_Interval}
                          \right) = 0$.
            \newline
            For example, if $\inptext{Local\_Limiter\_Interval}=1$, the local
            limiter will be applied at the end of every time step. The
            ${current\_iteration}$ is the total aggregate iteration
            number including previous jobs.
 %          \\
 %$< 0  $ & The limiter will NOT be applied even if \newline
 %          $\inptext{Local\_Limiter\_Option} \neq 0$
\end{tabular}
\end{nmldescbody}

%==========================================================================

\begin{nmldescitem}
Local_Limiter_Option = 0
\end{nmldescitem}
\inp*[nmlLimiter::LocalLimiterOption]{}[Local\_Limiter\_Option][]

\begin{nmldescbody}
GFR uses a three-pass filtering approach involving global, local, and
sponge filters.
%
Each of them and their sequence of application is described in
\sref{SEC:Filtering}.
%
\inptext{Local\_Limiter\_Option} specifies the type of limiter to apply to
locally marked cells.
\\
\begin{tabular}{ @{\quad} r @{ $=$ } p{0.85\linewidth} @{} }
  $   0$ & Do not apply any kind of limiter. \\
  $   1$ & Apply filtering. \\
  $   2$ & Apply scaled filtering.
           The filter value is scaled by the magnitude of the scaled marker. \\
  $   3$ & Apply the multi-dimensional limiting process (MLP). \\
  $9001$ & Apply recursive projection.
           For any troubled cell that is marked, project the solution down one
           order; recheck that cell to see if it is still troubled; keep
           repeating this process until the cell is no longer marked or has been
           reduced to a \PO[1] solution.
           \\
  $9002$ & Apply scaled projection.
           For any troubled cell that is marked,
           project the solution down to a order that is dynamically
           chosen by the value of the local marker.
           Cells with higher marker values will be projected to lower orders,
           whereas cells with lower marker values will be projected down
           only an order or two from the solution order.
\end{tabular}

\noindent
\NOTE
%
The locally marked cells are determined from the inputs for
  \inp[nmlLimiter::MarkerOption]{Marker\_Option} and
  \inp[nmlLimiter::MarkerVariable]{Marker\_Variable}.
%
The local filter settings are defined by input values for
  \inp[nmlLimiter::LocalFilterOrder]{Local\_Filter\_Order},
  \inp[nmlLimiter::LocalFilterAlpha]{Local\_Filter\_Alpha}, and
  \inp[nmlLimiter::LocalFilterEtac]{Local\_Filter\_Etac}.

\end{nmldescbody}

%==========================================================================

\begin{nmldescitem}
Local_Filter_Order  = 16
Local_Filter_Alpha  = 36.0
Local_Filter_Etac   =  0.0
\end{nmldescitem}
\inp*[nmlLimiter::LocalFilterOrder]{}[Local\_Filter\_Order][]
\inp*[nmlLimiter::LocalFilterAlpha]{}[Local\_Filter\_Alpha][]
\inp*[nmlLimiter::LocalFilterEtac]{}[Local\_Filter\_Etac][]

\begin{nmldescbody}
As discussed in \sref{SEC:Filtering}, the local filter is applied
to those cells identified by the troubled cell marker.
%
A good approach is to use the local options to apply a stronger filter
to only those cells that have been marked as a troubled.
%
This will help prevent the oscillatory solutions in these cells
from causing the simulation to diverge.
%
Then, if necessary, a weaker global filter can be applied
to all cells to damp any other disturbances.

The details for these inputs are the same as for the similarly
named \inptext{Global\_Filter} options, but one likely wants
to apply a stronger local filter setting.
%
Use
\inptext{Local\_Filter\_Order}=16 and halve if necessary, and
\inptext{Local\_Filter\_Alpha}=36 and halve if possible.
For \inptext{Local\_Filter\_Etac}, use the same recommendations as
for \inp[nmlLimiter::GlobalFilterEtac]{Global\_Filter\_Etac}.

\IDEA{
Create two paragraphs here: one with recommendations for
$\inptext{Local\_Limiter\_Option}=1$, and one with recommendations for 
$\inptext{Local\_Limiter\_Option}=2$.
}
\end{nmldescbody}

%==========================================================================

\FINALFIX{\newpage}
\begin{nmldescitem}
Marker_Stored_Interval = 1
\end{nmldescitem}
\inp*[nmlLimiter::MarkerStoredInterval]{}[Marker\_Stored\_Interval][]

\begin{nmldescbody}
This determines if the marker values written to the CGNS solution files
are from the final Runge-Kutta stage or from each Runge-Kutta stage.
\\
\begin{tabular}{ @{\qquad} r @{ \quad } p{0.85\linewidth} @{} }
  $<0$ & Store markers at the same frequency of
         \inp[nmlLimiter::LocalLimiterInterval]{Local\_Limiter\_Interval}.
         This can be used to synchronize the marker storage and
         limiter application.
         \\
  $=0$ & Store markers after each Runge-Kutta stage. \\
  $>0$ & Store markers after the last Runge-Kutta stage.
\end{tabular}
\\
\NOTE
Storing markers each stage increases memory and file size
(if \inp[nmlOutputVariables::MarkerVariables]{Marker\_Variables}
are being output).
\end{nmldescbody}

%==========================================================================

%##%\FINALFIX{\newpage}
\begin{nmldescitem}
Marker_Option = 0
\end{nmldescitem}
\inp*[nmlLimiter::MarkerOption]{}[Marker\_Option][]

\begin{nmldescbody}
Specifies which type of local marker to use.
\\
\begin{tabular}{ @{\qquad} r @{ = } p{0.85\linewidth} @{} }
  0 & Use the default marker option for the chosen
      \inp[nmlLimiter::LocalLimiterOption]{Local\_Limiter\_Option}. \\
  1 & Use the resolution marker described in \sref{SEC:ResolutionMarker}.\\
  2 & Use the jump marker described in \sref{SEC:JumpMarker}.\\
  3 & Use the MLP marker described in \sref{SEC:MLPMarker}.\\
  4 & Use the Ducros marker described in \sref{SEC:DucrosMarker}.
\end{tabular}\\
%
\Tref{TAB:Permissible_Marker_Option} lists the
permissible and default values for \inptext{Marker\_Option},
depending on
\\ % Prevent next word from entering right margin.
\inptext{Local\_Limiter\_Option}.

\IDEA{Add Ducros and Ducros$\times$Jump markers.}

\begin{table}[h]
\begin{center}
\caption{Permissible values for \inptext{Marker\_Option}.}
\label{TAB:Permissible_Marker_Option}
\begin{tabular}{rlc}
  \multicolumn{2}{c}{\inptext{Local\_Limiter\_Option}} &
  \inptext{Marker\_Option} \\
  \hline
     0 & None                & $1^*,2,3,4$ \\
     1 & Filtering           & $1^*,2,3,4$ \\
     2 & Scaled Filtering    & $1^*,2,3,4$ \\
     3 & MLP                 & $3^*$     \\
  9001 & Recusive Projection & $1^*,2$   \\
  9002 & Scaled   Projection & $1^*,2$   \\
  \hline
\end{tabular}
\\
${}^*$ indicates default value.
\end{center}
\end{table}

\end{nmldescbody}

%==========================================================================

\begin{nmldescitem}
Marker_Variable = 0
\end{nmldescitem}
\inp*[nmlLimiter::MarkerVariable]{}[Marker\_Variable][]

\begin{nmldescbody}
This determines which variable is used to evaluate the local marker.
\\
\begin{tabular}{ @{\qquad} r @{ = } p{0.85\linewidth} @{} }
  -2 & Use the Ducros variable based upon velocity divergence (only valid for
       $\inptext{Marker\_Option}=4$).\\
  -1 & Use pressure.\\
   0 & Use the default marker variable for the chosen
       \inp[nmlLimiter::MarkerOption]{Marker\_Option}.
       \\
   1 & Use density.\\
   2 & Use $x$-momentum.\\
   3 & Use $y$-momentum.\\
   4 & Use $z$-momentum.\\
   5 & Use total energy.\\
   6 & Use the max value after evaluating the marker
       for each conserved variable.
\end{tabular}\\
%
\Tref{TAB:Permissible_Marker_Variable} lists the
permissible and default values for \inptext{Marker\_Variable},
depending on
\\ % Prevent next word from entering right margin.
\inptext{Marker\_Option}.

\begin{table}[h]
\begin{center}
\caption{Permissible values for \inptext{Marker\_Variable}.}
\label{TAB:Permissible_Marker_Variable}
\begin{tabular}{rlc}
  \multicolumn{2}{c}{\inptext{Marker\_Option}} & \inptext{Marker\_Variable} \\
  \hline
     1 & Resolution & $ 1^*,2,3,4,5,6$ \\
     2 & Jump       & $-1^*,1,2,3,4,5$ \\
     3 & MLP        & $ 1^*,2,3,4,5$ \\
     4 & Durcos     & $-2^*$ \\
  \hline
\end{tabular}
\\
${}^*$ indicates default value.
\end{center}
\end{table}

\NOTE
Using all conserved variables (option 6) will result in marking
the most number of cells.
This will increase stability, but has the tendency to add unnecessary
dissipation where it is not needed.
The default of using only density seems to result in the best mix of
stability with minimal unnecessary dissipation.
\end{nmldescbody}

%==========================================================================

\begin{nmldescitem}
Marker_Min = 1.0
\end{nmldescitem}
\inp*[nmlLimiter::MarkerMin]{}[Marker\_Min][]

\begin{nmldescbody}
This gives the value for $\Psi_{min} = \Psi_0 - \Delta \Psi$,
which is the bottom-end of the $s$-curve used to scale
the marker value from the interval ($-\infty$,0] to the interval [0,1].
%
Any cell with a marker value less than the value of this input will have
a scaled marker value of 0, which means that cell will receive no limiting.
%
See the discussion of the scaled marker in \sref{SEC:ScaledMarker}
for a more information
\end{nmldescbody}

%==========================================================================

\begin{nmldescitem}
Marker_Max = 1.0
\end{nmldescitem}
\inp*[nmlLimiter::MarkerMax]{}[Marker\_Max][]

\begin{nmldescbody}
This gives the value for $\Psi_{max} = \Psi_0 + \Delta \Psi$,
which is the top-end of the $s$-curve used to scale
the marker value from the interval ($-\infty$,0] to the interval [0,1].
%
Any cell with a marker value greater than the value of this input will have
a scaled marker value of 1, which means that cell will receive
the maximum strength of the limiting method being employed.
%
See the discussion of the scaled marker in \sref{SEC:ScaledMarker}
for a more information
\end{nmldescbody}

%==========================================================================

\begin{nmldescitem}
Marker_Dpsi = -1.0
\end{nmldescitem}
\inp*[nmlLimiter::MarkerDpsi]{}[Marker\_Dpsi][]

\begin{nmldescbody}
This gives the value for the constant $\Delta\Psi$ that is used to scale
the marker value from the interval ($-\infty$,0] to the interval [0,1].
%
See the discussion of the scaled marker in \sref{SEC:ScaledMarker}
for a more information
\end{nmldescbody}

%==========================================================================

\FINALFIX{\newpage}
\begin{nmldescitem}
Thermofix%Method          = 0
Thermofix%Iterations      = 10
Thermofix%Check_Cells     = T
Thermofix%Check_Faces     = T
Thermofix%Check_Nodes     = T
Thermofix%Verbose         = T
Thermofix%Scale_Factor    = 0.5
Thermofix%Min_Pressure    = 10000.  ! [Pa]
Thermofix%Min_Temperature = 100.    ! [K]
Thermofix%Min_Density     = 0.001   ! [kg/m^3]
\end{nmldescitem}
\inp*[nmlLimiter::Thermofix/Method]{}[Thermofix\%Method][]
\inp*[nmlLimiter::Thermofix/Iterations]{}[Thermofix\%Iterations][]
\inp*[nmlLimiter::Thermofix/CheckCells]{}[Thermofix\%Check\_Cells][]
\inp*[nmlLimiter::Thermofix/CheckFaces]{}[Thermofix\%Check\_Faces][]
\inp*[nmlLimiter::Thermofix/CheckNodes]{}[Thermofix\%Check\_Nodes][]
\inp*[nmlLimiter::Thermofix/Verbose]{}[Thermofix\%Verbose][]
\inp*[nmlLimiter::Thermofix/ScaleFactor]{}[Thermofix\%Scale\_Factor][]
\inp*[nmlLimiter::Thermofix/MinPressure]{}[Thermofix\%MinPressure][]
\inp*[nmlLimiter::Thermofix/MinTemperature]{}[Thermofix\%Min\_Temperature][]
\inp*[nmlLimiter::Thermofix/MinDensity]{}[Thermofix\%Min\_Density][]

\begin{nmldescbody}
Sometimes during the course of a simulation, the solution will
encounter a negative thermodynamic property
(pressure, temperature, or density),
which can cause a NaN (\forexample when computing speed of sound).
%
The \inptext{Thermofix} inputs are intended to help remedy this
situation.
%
While they provide additional stability when needed, they also introduce
additional dissipation, and should be avoided if possible.

\inptext{Method} specifies the method used to repair
cells with bad thermodynamic properties.
Checks are performed during each Runge-Kutta stage.
\\
\begin{tabular}{ @{\qquad} r @{ = } p{0.85\linewidth} @{} }
  0 & None \\
  1 & Limit the high-order modes using the exponential sponge filter. \\
  2 & Scale the high-order modes by \inptext{Scale\_Factor}.
\end{tabular}

\inptext{Iterations} specifies how many passes to
make before abandoning further repair efforts.

\inptext{Scale\_Factor} specifies the scaling factor to apply
to the high-order modes when using \inptext{Method}=2.

\inptext{Check\_Cells},
\inptext{Check\_Faces}, and
\inptext{Check\_Nodes}
specifies whether to check the thermodynamic properties at points within the
cell volume, the cell faces, or the cell nodes (corners).
%
If any of these checks identifies a troubled point,
the solution within the cell volume is treated and
re-interpolated to the faces and nodes.

\inptext{Verbose} is used to output information about
the bad cells to the display output.

\inptext{Min\_Pressure},
\inptext{Min\_Temperature}, and
\inptext{Min\_Density}
are used to specify threshold values below which action should be taken.
%
These limits are imposed individually, so the three specified values need not
satisfy the equation of state.
%
Using small yet positive values is intended to identify problem
regions before the solution has deteriorated too severely.
%
The actual values for these parameter may need to be adjusted
for each flow problem.
%
These values should be small enough to not be physically encountered within
the flow but large enough that the solution can still be recovered before
fully diverging.

The thermofix limiter is applied at the end of each Runge-Kutta stage.
%
The order in which the thermofix limiter is applied relative to
other limiting\slash filtering methods can be found in the enumerated list
towards the beginning \sref{SEC:Numerics}.

\end{nmldescbody}

%==========================================================================

\FINALFIX{\newpage}
\begin{nmldescitem}
Positivity%Option         = 0
Positivity%Density_Limit  = 1.0e-6 ! [kg/m^3]
Positivity%Pressure_Limit = 1.0e-6 ! [Pa]
\end{nmldescitem}
\inp*[nmlLimiter::Positivity/Option]{}[Positivity\%Option][]
\inp*[nmlLimiter::Positivity/DensityLimit]{}[Positivity\%Density\_Limit][]
\inp*[nmlLimiter::Positivity/PressureLimit]{}[Positivity\%Pressure\_Limit][]

\begin{nmldescbody}
 
%ADD REFERENCE TO ZHANG AND SHU, and JOHNSON AND CHING.

Sometimes during the course of a simulation, the solution will encounter a
negative thermodynamic property (pressure, temperature, or density), which can
cause a NaN (\forexample when computing speed of sound).
%
The \inptext{Positivity} inputs are intended to help remedy this situation using
a positivity-preserving limiter.
%
While this method provides additional stability when needed, it will also
introduce additional dissipation, and should therefore be avoided if possible.

\inptext{Option} simply specifies if the positivity limiter is activated.
\\
\begin{tabular}{ @{\qquad} r @{ = } p{0.85\linewidth} @{} }
  0 & Off (Default) \\
  1 & On
\end{tabular}

\inptext{Density\_Limit} specifies the minimum dimensional density allowable for
any of the solution or flux points of a cell.
 
\inptext{Pressure\_Limit} specifies the minimum dimensional pressure allowable
for any of the solution or flux points of a cell.

When the positivity limiter is activated, if any solution or flux point of a
cell has a density or pressure value below their respective limits, the
positivity limiter will be applied to this cell.

The positivity limiter is applied at the end of each Runge-Kutta stage.
%
The order in which the positivity-preserving limiter is applied relative to
other limiting\slash filtering methods can be found in the enumerated list
towards the beginning \sref{SEC:Numerics}.
%
%The positivity limiter will be applied \emph{after} any of the following (if
%also active): sponge filter or weak global filter (\thatis
%$\inptext{Global\_Filter\_Alpha} \leq 1$).
%%
%The positivity limiter will be applied \emph{before} any of the following (if
%also active): entropy limiter, local limiter (\thatis
%$\inptext{Local\_Limiter\_Option}\neq 0$), strong global filter (\thatis
%$\inptext{Global\_Filter\_Alpha}>1$), or \inptext{Thermofix}.

When the positivity limiter is applied to a cell, the solution in the cell is
modified in the following way:
%
\begin{equation}
  \tilde{u} = \check{u} + \theta_{PL}\left(u-\check{u}\right)
  \label{EQN:PLSolution}
\end{equation}
%
where $u$ is the conserved variable being limited, $\check{u}$ is the average
value of the polynomial $u$ in the cell, $\theta_{PL}$ is the limiter scaling
factor, and $\tilde{u}$ is the resulting limited polynomial.
%
The limiter scaling factor, $\theta_{PL}$, is computed as
%
\begin{equation}
  \theta_{PL} = 0.99 \left(
                           \frac{\check{a}-a_{\textrm{Lim}}}
                                {\check{a}-a_{\textrm{min}}}
                     \right)
  \label{EQN:PLScalingFactor}
\end{equation}
%
where $a$ is either density or pressure depending on which positivity limit was
violated in the cell, $\check{a}$ is the average value of the variable $a$ in
the cell, $a_{\textrm{Lim}}$ is the value of the respective input
\inptext{Density\_Limit} or \inptext{Pressure\_Limit}, and $a_{\textrm{min}}$ is
the minimum value of the respective variable from all solution and flux points
of the cell.

\NOTE
If \inptext{Density\_Limit} is violated in the cell, only the density polynomial
is limited by the positivity limiter, with the scaling factor $\theta_{PL}$
evaluated with $a=\rho$.
%
If \inptext{Pressure\_Limit} is violated in the cell, all conserved variables
are limited by the positivity limiter, with the scaling factor $\theta_{PL}$
evaluated with $a=p$.

\NOTE
The density and pressure limits needed for the positivity limiter are generally
much smaller than what is needed for \inptext{Thermofix}.

\end{nmldescbody}

%==========================================================================

\FINALFIX{\newpage}
\begin{nmldescitem}
Entropy_Limiter%Option      = 0
Entropy_Limiter%Interpolant = 0
Entropy_Limiter%Tolerance   = 0.001
Entropy_Limiter%Check_Faces = F
\end{nmldescitem}
\inp*[nmlLimiter::EntropyLimiter/Option]{}[Entropy\_Limiter\%Option][]
\inp*[nmlLimiter::EntropyLimiter/Interpolant]{}[Entropy\_Limiter\%Interpolant][]
\inp*[nmlLimiter::EntropyLimiter/Tolerance]{}[Entropy\_Limiter\%Tolerance][]
\inp*[nmlLimiter::EntropyLimiter/CheckFaces]{}[Entropy\_Limiter\%Check\_Faces][]

\begin{nmldescbody}
 
%ADD REFERENCE TO ZHANG AND SHU, and JOHNSON AND CHING.

Sometimes during the course of a simulation, the solution will encounter a
negative thermodynamic property (pressure, temperature, or density), which can
cause a NaN (\forexample when computing speed of sound).
%
The \inptext{Entropy\_Limiter} inputs are intended to help remedy this situation
using a limiter to prevent violations of the local entropy bound in a cell.
%
While this method provides additional stability when needed, it will also
introduce additional dissipation, and should therefore be avoided if possible.

\inptext{Option} specifies if the entropy limiter is active, and if so, which
method to use to limit any cells that violate the minimum entropy bounds.
\\
\begin{tabular}{ @{\qquad} r @{ = } p{0.85\linewidth} @{} }
  0 & Off (Default) \\
  1 & Flat filter \\
  2 & Scaled filter
\end{tabular}

\inptext{Interpolant} specifies the interpolant used by the entropy limiter.
%
For each cell that is found to be in violation of the local entropy bounds, the
interpolant is used to compute the filter strength, $\theta_{EL}$, necessary to
prevent the cell from committing the violation.
\\
\begin{tabular}{ @{\qquad} r @{ = } p{0.85\linewidth} @{} }
  0 & method used by Lv-Ihme; equation 44 from \linecite{Lv_2015} (Default) \\
  1 & linear \\
  2 & density weighted \\
  3 & log scaling \\
  4 & Pad\'{e} approximation
\end{tabular}
   
\inptext{Tolerance} specifies how much the local entropy bounds can be violated
before limiting the cell.
%
As this value gets closer to zero, the limiter gets more strict about not
allowing the solution to violate the bounds by adding more dissipation to
prevent the violation from occurring.

\inptext{Check\_Faces} specifies if the entropy values from the faces of each
cell (both internal and from the adjacent neighboring cells) are included when
computing the minimum entropy for said cell.
%
Including the face entropy values in the minimum entropy will help with
maintaining a stable solution, but at the expense of additional dissipation.

%%%%%%%%%%%%%%%%%%%%%%%%%%%%%%%%%%%%%%%%%%%%%%%%%%%%%%%%%%%%%%%%%%%%%%%%%%%%%%%%

When the entropy limiter uses the flat filter to modify the solution in a cell,
the limited solution polynomial is given by
%
\begin{equation}
  \tilde{u} = \check{u} + \theta_{EL}\left(u-\check{u}\right)
  \label{EQN:ELFlatSolution}
\end{equation}
%
where $u$ is the conserved variable being limited, $\check{u}$ is the average
value of the polynomial $u$ in the cell, $\tilde{u}$ is the resulting
limited polynomial, and the filter strength is in the range $\theta_{EL} =
\left[0,1\right]$.
%
When using the scaled filter, the entropy limiter modifies the solution in one
of two ways depending on the value of $\theta_{EL}$:
%
\begin{equation}
  \tilde{u} = 
  \begin{cases}
    u + \theta_{EL}\left({\cal F}\left(u\right) - u\right) & 
       \textrm{if } \theta_{EL} \leq 1 \\
    \check{u} + \left(2-\theta_{EL}\right) \left({\cal F}\left(u\right) -
    \check{u}\right) &
       \textrm{if } \theta_{EL} > 1
  \end{cases}
  \label{EQN:ELScaledSolution}
\end{equation}
%
where ${\cal F}\left(u\right)$ is the result of applying the full local filter
to the unmodified solution $u$, and the filter strength is in the range
$\theta_{EL} = \left[0,2\right]$.
  
%%%%%%%%%%%%%%%%%%%%%%%%%%%%%%%%%%%%%%%%%%%%%%%%%%%%%%%%%%%%%%%%%%%%%%%%%%%%%%%%

The entropy limiter is applied at the end of each Runge-Kutta stage.
%
The order in which the entropy limiter is applied relative to other
limiting\slash filtering methods can be found in the enumerated list towards the
beginning \sref{SEC:Numerics}.
%
%The entropy limiter will be applied \emph{after} any of the following (if also
%active): sponge filter, weak global filter (\thatis
%$\inptext{Global\_Filter\_Alpha} \leq 1$), or positivity-preserving limiter.
%%
%The entropy limiter will be applied \emph{before} any of the following (if also
%active): local limiter (\thatis $\inptext{Local\_Limiter\_Option}\neq 0$),
%strong global filter (\thatis $\inptext{Global\_Filter\_Alpha}>1$), or
%\inptext{Thermofix}.

%%%%%%%%%%%%%%%%%%%%%%%%%%%%%%%%%%%%%%%%%%%%%%%%%%%%%%%%%%%%%%%%%%%%%%%%%%%%%%%%

\NOTE
When using $\inptext{Option}=2$, the local filter will be used to compute a
filtered solution to use as part of the limiting process; therefore, the input
coefficients for the local filter (\inptext{Local\_Filter\_Order},
\inptext{Local\_Filter\_Alpha}, and \inptext{Local\_Filter\_Etac}) will become
an important part of the entropy limiter settings and need to be adjusted
accordingly.

\NOTE
When the entropy limiter is applied to a cell, all conserved variables are
modified together with the same filter strength.

\NOTE
So far, our experience seems to indicate that the \inptext{Interpolant} has
little effect on the stability or accuracy of the solution; however, this could
be something that is problem dependent and we have yet to find a case that
demonstrates the effect of different interpolant methods.
%
Because of this, it is recommended to first look at adjusting other inputs when
encountering stability issues.

\end{nmldescbody}

%==========================================================================

